% !PRINTED 454
\weberpage{454}

%<<CONT
Doppelpunkt, Wendepunkt, Rückkehrpunkt, Berührungspunkt einer Doppeltangente u.\,s.\,w. ist.

Unter den Formen $\varphi$ sind uns nun vor Allem die von Wichtigkeit, die bei der ganzen Gruppe ungeändert bleiben, die Invarianten, deren es hier, da die Gruppe einfach ist, nur absolute giebt (\S.\,40). Ist $\Phi\,(x_1, x_2, x_3)$ eine solche Form, so soll die durch die Gleichung $\Phi = 0$ dargestellte Curve eine invariante Curve der Gruppe heissen. Es giebt 168 Abbildungen der Ebene auf sich selbst, bei denen alle Punkte der Curve in Punkte derselben Curve übergehen. Liegt irgend ein Punkt auf dieser Curve, so liegen auch alle seine Bildpunkte darauf.

Im Allgemeinen ist die Zahl der so mit einander verbundenen Punkte der Curve 168, jedenfalls nicht grösser. Ist sie kleiner, so müssen die Punkte \gesperrt{Pole} sein, und die Anzahl der Bildpunkte ist ein Theiler von 168 (nämlich 24 für die $P_7$, 21 für die $P_8$, 28 für die $P_6$, 42 für die $P_4$, 56 für die $P_3$ und 84 für ein System zusammengehöriger $P_2$). Wenn ein Pol von einer dieser Arten auf der Curve liegt, so liegen alle Pole von derselben Art darauf.

\begin{center}
\S.\,122.\\
\gesperrt{Die erste Invariante der Gruppe $G_{168}$ und die Grundcurve.}
\end{center}

Um nun die Invarianten unserer Gruppe zu bilden, suchen wir Formen der drei Variablen $x_1, x_2, x_3$ auf, die durch Anwendung der drei Substitutionen $\chi, \tau, \omega$ (\S.\,115) ungeändert bleiben. Dies genügt, da die ganze Gruppe sich aus diesen drei Substitutionen zusammensetzen lässt (\S.\,72). Nun ist $\chi$ eine cyklische Vertauschung der drei Variablen $x_1, x_2, x_3$, und $\tau$ bedeutet die Substitution

\[
\begin{pmatrix} x_1, & x_2, & x_3 \\ \varepsilon x_1, & \varepsilon^2 x_2, & \varepsilon^4 x_3 \end{pmatrix}.
\]

Wenn also in einer Invariante ein Glied $x_1^{h_1} x_2^{h_2} x_3^{h_3}$ vorkommt, worin $h_1, h_2, h_3$ ganze nicht negative Zahlen sind, so verlangt die Unveränderlichkeit durch $\tau$, dass

\weq{1}{h_1 + 2h_2 + 4h_3 \equiv 0\ (\mathrm{mod}\ 7)}
%CONT>>
